\section*{Bab \ref{ch:b2:speciation} --- Spesiasi dan Makroevolusi}

\begin{solution}{exo:b2:speciation:1}
Sebuah spesies adalah sekelompok populasi alami yang saling mengawini
yang terisolasi secara reproduktif dari kelompok lain semacam itu. Ia
gagal bagi garis keturunan aseksual (bakteri, rotifera bdeloid,
randa tapak: konsep morfologis atau filogenetik, yaitu gugus bentuk
yang serupa atau garis keturunan yang dapat didiagnosis); bagi fosil
(konsep morfologis, karena tidak ada persilangan yang dapat dibuat);
dan bagi populasi alopatrik yang tidak pernah berjumpa (konsep ekologis
atau filogenetik, karena perkawinan silangnya tidak dapat diuji) ---
dan ia dikaburkan pohon ek yang berhibrida dan spesies cincin.
\end{solution}

\begin{solution}{exo:b2:speciation:2}
Katak pada bulan yang berbeda: prazigotik (waktu). Bagal: pascazigotik
(kemandulan hibrida). Sperma yang tidak dapat mengikat: prazigotik
(gametik). Hibrida yang gagah dengan \omterm{def:b2:angiosperm-reproduction:seed}{biji} yang mandul: pascazigotik
(kemandulan hibrida, atau keruntuhan hibrida jika kegagalannya berada
pada generasi berikutnya). Jangkrik yang nyanyiannya berbeda:
prazigotik (perilaku).
\end{solution}

\begin{solution}{exo:b2:speciation:3}
Alopatrik: sebuah penghalang membagi sebuah sebaran (udang pistol pada
kedua sisi Tanah Genting Panama). Peripatrik: beberapa pendiri di balik
sebarannya (burung pipit Darwin dari seekor moyang daratan). Simpatrik:
tanpa penghalang (lalat belatung apel pada hawthorn dan apel; sebuah
gandum \omterm{def:b2:genome-diversification:polyploidy}{alopoliploid}). Parapatrik: sepanjang sebuah landaian dengan
sebuah \omterm{prop:b2:speciation:reinforcement}{zona hibrida} di tengahnya (rumput di atas dan di luar timbunan
tambang, yang berbunga pada waktu yang berbeda).
\end{solution}

\begin{solution}{exo:b2:speciation:4}
Spesies 1 selalu menang; spesies 2 selalu menang; \omterm{thm:b2:speciation:lv}{hidup berdampingan}
yang mantap; salah satunya menang menurut jumlah awalnya. \omterm{thm:b2:speciation:lv}{Hidup
berdampingan} mantap ketika tiap spesiesnya dapat menyerbu spesies
lainnya pada daya dukungnya, yakni ketika tiap spesiesnya membatasi
dirinya sendiri lebih daripada ia membatasi yang lain ($K_1 > \alpha K_2$ dan $K_2 > \beta K_1$).
\end{solution}

\begin{solution}{exo:b2:speciation:5}
$\alpha K_2 = 400 < K_1 = 1000$: spesies 1 menyerbu; $\beta K_1 = 600
< K_2 = 800$: spesies 2 menyerbu. $N_1^{*} = (1000 - 400)/(1 - 0.3) =
857$, $N_2^{*} = (800 - 600)/0.7 = 286$; pada kesetimbangannya spesies
1 berada di bawah $K$-nya sebesar $\alpha N_2^{*} = 143$ dan spesies 2
sebesar $\beta
N_1^{*} = 514$.
\end{solution}

\begin{solution}{exo:b2:speciation:6}
$\alpha K_2 = 1200 > K_1$: spesies 1 tidak dapat menyerbu spesies 2
pada daya dukungnya; $\beta K_1 = 300 < K_2$: spesies 2 dapat menyerbu
spesies 1. Spesies 2 selalu menang, karena individunya melukai spesies
1 lebih daripada individu spesies 1 sendiri sementara ia sendiri nyaris
tidak terpengaruh. Dengan $\alpha = \beta = 1.5$: $\alpha K_2 = 1200 > 1000$
dan $\beta K_1 = 1500 > 800$: tak satu pun dapat menyerbu yang lain,
masing-masing menyingkirkan yang lain ketika ia mapan, dan pendirinya
yang menang.
\end{solution}

\begin{solution}{exo:b2:speciation:7}
Tetraploid $4n$ membuat gamet $2n$; jika disilangkan dengan gamet $n$
\omterm{def:b2:meiosis-heredity:meiosis}{diploidnya} ia memberikan triploid $3n$, yang \omterm{def:b2:meiosis-heredity:meiosis}{meiosisnya} tidak dapat
memasangkan tiga perangkat lalu menghasilkan gamet aneuploid yang
kebanyakan tidak dapat hidup. Tetraploidnya hanya dapat berbiak dengan
dirinya sendiri (atau dengan \omterm{def:b2:angiosperm-reproduction:pollination}{penyerbukan} sendiri): terisolasi dalam
satu langkah. Ia biasanya gagal karena ia sendirian --- satu-satunya
calon pasangannya adalah \omterm{def:b2:meiosis-heredity:meiosis}{diploid} dan tiap persilangan semacam itu
terbuang --- lalu tenggelam; ia bertahan ketika ia menyerbuki diri,
berkembang biak secara vegetatif, atau muncul berulang kali.
\end{solution}

\begin{solution}{exo:b2:speciation:8}
$0.01\times 10\times 10 = 1$ ketaksesuaian; dengan 30 masing-masing,
$0.01\times
900 = 9$. Pengisolasiannya bertambah sebagai kuadrat waktu
pemencarannya: dua populasi yang memencar selama tiga kali lebih lama
sembilan kali lebih terisolasi --- yaitu bola saljunya.
\end{solution}

\begin{solution}{exo:b2:speciation:9}
$w = 2\sqrt{8/0.05} = \qty{25}{km}$; $w = 0.3\sqrt{8/0.4} =
\qty{1.3}{km}$. Penguatannya lebih mungkin pada yang kedua: hibridanya
sangat tidak bugar, sehingga seekor betina yang menolak jenis yang lain
memperoleh banyak, dan zonanya cukup sempit sehingga kedua bentuknya
cukup sering berjumpa agar penolakan itu terseleksi.
\end{solution}

\begin{solution}{exo:b2:speciation:10}
$\alpha K_2 = 450 < 500$, sehingga masing-masing menyerbu yang lain:
keduanya \omterm{thm:b2:speciation:lv}{hidup berdampingan}, pada $N^{*} = (500 - 450)/(1 - 0.81) = 263$ masing-masing --- nyaris saja, dekat tepi penyingkirannya, dan
sebuah perubahan kecil pada $K$ akan menjungkitkannya. Dengan $\alpha = \beta = 0.4$: $N^{*} = (500 - 200)/(1 - 0.16) = 357$ masing-masing.
Pergeserannya memutar tiap isoklinnya menjauhi isoklin yang lain
($K/\alpha$ bergerak ke luar), sehingga perpotongannya bergerak menjauhi
sumbunya dan tiap spesiesnya duduk lebih dekat kepada $K$-nya sendiri:
sebuah \omterm{thm:b2:speciation:lv}{hidup berdampingan} yang lebih kukuh dengan lebih banyak
masing-masingnya.
\end{solution}

\begin{solution}{exo:b2:speciation:11}
$R = h^{2}S = 0.7\times(-1) = \qty{-0.7}{mm}$: seperti yang teramati.
Pada kekeringan tahun 1977, tanpa kehadiran burung pipit tanah besar,
burung pipit sedang yang sama terseleksi bagi paruh yang \emph{lebih
besar}, karena \omterm{def:b2:angiosperm-reproduction:seed}{biji} besarnya adalah makanan yang tersisa; pada tahun
2004 \omterm{def:b2:angiosperm-reproduction:seed}{biji} besarnya diambil penyerbunya, dan burung pipit sedang yang
bersaing memperebutkannya kalah: arah seleksinya ditetapkan pesaingnya.
\end{solution}

\begin{solution}{exo:b2:speciation:12}
Dibuat secara kebetulan: model dua lokusnya memperlihatkan
pengisolasiannya muncul sebagai hasil sampingan penggantian yang
terpancang karena alasan lain, atau lewat hanyutan, dan geografinya ---
sebuah penghalang, sebuah pendirian --- memasok pemisahan yang
membiarkan kebetulannya menimbun. Dijaga oleh seleksi: di tempat kedua
bentuknya berjumpa, seleksi terhadap hibrida yang tidak bugar membangun
penghalang prazigotik (penguatan), dan persaingannya menggeser sifatnya
sehingga keduanya dapat \omterm{thm:b2:speciation:lv}{hidup berdampingan}. Semboyannya melewatkan
\omterm{prop:b2:speciation:modes}{spesiasi simpatrik} lewat seleksi pemecah, yang padanya seleksi membuat
sekaligus menjaga spesiesnya, dan \omterm{def:b2:genome-diversification:polyploidy}{poliploidi}, yang padanya kebetulannya
adalah keseluruhan ceritanya.
\end{solution}

\begin{solution}{pb:b2:speciation:1}
\textbf{1.} Spesies 1: $N_1 + 0.9N_2 = 1200$, dengan potongan $N_1 = 1200$
dan $N_2 = 1333$. Spesies 2: $N_2 + 0.3N_1 = 400$, dengan potongan $N_2 =
400$ dan $N_1 = 1333$.
\textbf{2.} $\beta K_1 = 360 < K_2 = 400$: spesies 2 menyerbu;
$\alpha K_2 = 360 < K_1 = 1200$: spesies 1 menyerbu. \omterm{thm:b2:speciation:lv}{Hidup
berdampingan} yang mantap.
\textbf{3.} $N_1^{*} = (1200 - 360)/(1 - 0.27) = 1151$; $N_2^{*} =
(400 - 360)/0.73 = 55$.
\textbf{4.} $\mathrm{d}N_1/\mathrm{d}t = 0.5\times 1200\,(1 - 1218/1200)
= -9$ tiap tahun: penduduknya, yang berada pada daya dukungnya,
didorong menurun oleh pendatangnya.
\textbf{5.} $K_1 = 600$, $K_2 = 200$: $N_1^{*} = (600 - 180)/0.73 =
575$, $N_2^{*} = (200 - 180)/0.73 = 27$: keduanya menyetengah, dan
penyerbunya, pada 27 ekor burung, dekat dengan kepunahan --- kekeringan
adalah saat persaingannya menggigit.
\textbf{6.} Penyerbunya adalah burung yang lebih besar dengan paruh
yang lebih besar: ia mengambil \omterm{def:b2:angiosperm-reproduction:seed}{biji} besar yang juga dipakai
penduduknya, dan mengambilnya lebih cepat, sementara \omterm{def:b2:angiosperm-reproduction:seed}{biji} kecil
penduduknya tidak banyak berguna baginya.
\textbf{7.} $S = 9.0 - 9.5 = \qty{-0.5}{mm}$; $R = 0.7\times(-0.5) =
\qty{-0.35}{mm}$.
\textbf{8.} $9.5 - 0.35 = \qty{9.15}{mm}$.
\textbf{9.} $N_1^{*} = (1200 - 200)/(1 - 0.15) = 1176$; $N_2^{*} =
(400 - 360)/0.85 = 47$.
\textbf{10.} Potongan $N_2$ isoklin spesies 1 naik dari
$K_1/0.9 = 1333$ menjadi $K_1/0.5 = 2400$: garisnya berputar menjauhi
garis penyerbunya, penduduknya kini kurang terpengaruh tiap penyerbunya,
dan kesetimbangannya bergerak menuju $K$ penduduknya.
\textbf{11.} $2/0.35 \approx 6$ tahun. Hal itu tidak terjadi karena
seleksi semacam ini hanya berlangsung pada tahun kekeringannya,
berbalik pada tahun yang basah ketika \omterm{def:b2:angiosperm-reproduction:seed}{biji} kecilnya berlimpah (atau
ketika penyerbunya langka), dan karena ragam aditifnya akan terkuras.
\textbf{12.} Dua juta tahun adalah kira-kira $10^{6}$ generasi; sebuah
pergeseran sebesar \qty{0.35}{mm} tiap generasi yang berkelanjutan
bahkan selama seperseribunya akan mengubah paruhnya sebesar ratusan
milimeter. Seleksinya kuat tetapi tidak tetap; perubahan berarah yang
berkelanjutan adalah kekecualian yang jarang, dan perubahan bersih
sepanjang waktu geologis adalah sisa yang kecil dari peristiwa besar
yang saling berlawanan.
\textbf{13.} Nyanyiannya, yang dipelajari dari ayahnya dan dipakai
betinanya dalam memilih pasangan, dengan bentuk paruhnya sendiri
sebagai sebuah isyarat. Hibrida yang berparuh menengah tidak menangani
\omterm{def:b2:angiosperm-reproduction:seed}{biji} besar maupun \omterm{def:b2:angiosperm-reproduction:seed}{biji} kecilnya dengan terbaik lalu kelaparan lebih
dahulu pada masa kekeringan.
\textbf{14.} $w = 0.5\sqrt{8/0.2} = \qty{3.2}{km}$.
\textbf{15.} Pulaunya lebih sempit daripada zonanya kelak: tidak ada
klina ruang yang dapat terbentuk; hibridisasinya, di tempat ia terjadi,
mempengaruhi seluruh populasinya, dan hanya perkawinan berpilih yang
menjaga keduanya tetap terpisah.
\textbf{16.} Betina yang hanya mengawin dengan jantan yang nyanyian dan
paruhnya sejenis meninggalkan lebih banyak cucu daripada betina yang
berhibrida; kesukaannya dan isyaratnya (nyanyian, paruh) menjadi lebih
khas di tempat kedua spesiesnya hidup bersama daripada di tempat
keduanya hidup sendiri.
\textbf{17.} $0.005\times 15\times 15 = 1.1$ ketaksesuaian: secara
rerata satu, dengan hibridanya barangkali berkurang tetapi tidak
mandul. Belum spesies menurut kriteria biologisnya; sedang dalam
perjalanan.
\textbf{18.} Biologis: belum (hibridanya nyaris bugar). Morfologis:
barangkali, jika paruh atau bulunya telah memencar secara kasatmata.
Filogenetik: ya, jika tiap populasinya memiliki sebuah perbedaan
diagnostik yang terpancang. Konsepnya berselisih tepat selama proses
spesiasinya, dan itulah yang membuatnya berguna.
\textbf{19.} $1/(5\times 10^{6}) = 2\times 10^{-7}$ tiap spesies tiap
tahun; bagi $10^{7}$ spesies, dua kepunahan tiap tahun.
\textbf{20.} 200 sampai 2000 spesies tiap tahun.
\textbf{21.} $0.75\times 10^{7}/10^{5} = 75$ spesies tiap tahun ---
lebih rendah daripada laju sekarang pada ujung bawahnya. Peristiwa
sekarang, jika berkelanjutan, adalah sebuah \omterm{prop:b2:speciation:tempo}{kepunahan massal} yang
berjalan lebih cepat daripada lima yang besar.
\textbf{22.} Relung yang dikosongkan yang kalah bebas dari pesaing,
sehingga tiap varian yang bertahan hidup yang dapat memakai salah
satunya disukai, dan yang bertahan hidup meragamkan diri ke dalamnya.
\textbf{23.} \omterm{prop:b2:speciation:tempo}{Kesetimbangan tersela}; \omterm{prop:b2:speciation:modes}{spesiasi alopatrik} (peripatrik) di
dalam sebuah populasi kecil yang terisolasi yang keturunannya lalu
menyebar ke dalam sebaran populasi utamanya.
\textbf{24.} $1.04^{n} = 2$: $n = \ln 2/\ln 1.04 = 18$ generasi.
Catatannya memperlihatkan jutaan tahun karena seleksinya berbalik,
\omterm{prop:b2:population-genetics:kinds}{seleksi pemantap} menguasai sebagian besar waktunya, dan perubahan
bersihnya adalah selisih kecil antara peristiwa besar yang berlawanan.
\textbf{25.} Sebelum pergeserannya: \omterm{thm:b2:speciation:lv}{hidup berdampingan} yang mantap pada
$N_1^{*} =
1151$, $N_2^{*} = 55$; sesudahnya: 1176 dan 47. Pergeseran
paruhnya \qty{-0.35}{mm} dalam satu generasi. Lebar \omterm{prop:b2:speciation:reinforcement}{zona hibridanya}
\qty{3.2}{km}.
\end{solution}
